% Figure: schematic of composition and inverse. Three sets and two
% maps; the composition runs around the triangle and the inverse
% runs back along the diagonal.
\begin{figure}[htbp]
\centering
\begin{tikzpicture}[
  set/.style={draw=engblack, thick, ellipse,
    minimum width=2.0cm, minimum height=2.8cm},
  el/.style={circle, fill=engblack, inner sep=1.3pt},
  arrf/.style={->, thick, draw=engblack!85},
  arrg/.style={->, thick, draw=engblue},
  arrcomp/.style={->, very thick, draw=engred, dashed},
  ann/.style={font=\footnotesize, align=center},
  lbl/.style={font=\small\itshape, align=center},
]

% Three sets across the page
\node[set, label=above:{$A$}] (A) at (0, 0) {};
\node[set, label=above:{$B$}] (B) at (4.5, 0) {};
\node[set, label=above:{$C$}] (C) at (9.0, 0) {};

% Three points in A
\node[el] (a1) at (0, 0.7) {};
\node[el] (a2) at (0, 0) {};
\node[el] (a3) at (0, -0.7) {};

% Three points in B
\node[el] (b1) at (4.5, 0.7) {};
\node[el] (b2) at (4.5, 0) {};
\node[el] (b3) at (4.5, -0.7) {};

% Three points in C
\node[el] (c1) at (9.0, 0.7) {};
\node[el] (c2) at (9.0, 0) {};
\node[el] (c3) at (9.0, -0.7) {};

% f maps A to B
\draw[arrf] (a1) -- (b1);
\draw[arrf] (a2) -- (b2);
\draw[arrf] (a3) -- (b3);

% g maps B to C
\draw[arrg] (b1) -- (c1);
\draw[arrg] (b2) -- (c2);
\draw[arrg] (b3) -- (c3);

% Labels for f and g
\node[lbl] at (2.25, 1.6) {$f: A \to B$};
\node[lbl] at (6.75, 1.6) {$g: B \to C$};

% g o f composition: long arrow underneath
\draw[arrcomp] (a1) to[out=-30, in=-150] (c1);
\node[lbl, draw=engred, fill=white, inner sep=2pt]
  at (4.5, -2.0) {$g \circ f: A \to C$};

% f-inverse (return arrow above)
\draw[arrf, dotted, draw=engblack!60] (b2) to[out=110, in=70] (a2);
\node[ann] at (2.25, 1.0) {$f^{-1}$};

\end{tikzpicture}
\caption[Composition and inverse as directed maps between sets]{The
composition $g \circ f$ sends an element of $A$ to its image in $B$
under $f$ and then to its image in $C$ under $g$. The dashed red
arrow is the composed map read as a single function from $A$ to $C$.
The dotted arrow above shows the inverse of $f$ from $B$ back to
$A$, which exists exactly when $f$ is bijective. Composition is
associative ($h \circ (g \circ f) = (h \circ g) \circ f$) and is
rarely commutative. The inverse of the composition reverses both the
order and the individual inversions: $(g \circ f)^{-1} = f^{-1}
\circ g^{-1}$.}
\label{fig:vol02:ch02:composition-inverse}
\end{figure}
